
% argument 1: slopes (e.g. {4,6})
% argument 2: x position of the bottom left corner
% argument 3: y position of the bottom left corner
% argument 4: x length

%%%%%
% Plot slopes according to N or Kappa
%%%%%

\makeatletter
\newcommand{\printslope}[4]{
    \tikzset{fixed point arithmetic}
    % get arguments
    % #2 = xpos : plus petite valeur de h (coin haut-gauche)
    % #3 = ypos : valeur de kappa en xpos (grande valeur, en haut)
    % #4 = width : ecart vers la droite (hprevlast - hlast > 0)
    \def\nero@printslope@orderlist{#1}
    \edef\nero@printslope@xpos{#2}
    \edef\nero@printslope@ypos{#3}
    \edef\nero@printslope@width{#4}
    % px = xpos + width (avant-derniere valeur de h)
    \pgfmathparse{\nero@printslope@xpos+\nero@printslope@width}
    \edef\nero@printslope@px{\pgfmathresult}
    \edef\nero@printslope@py{\nero@printslope@ypos}
    \edef\nero@printslope@qx{\nero@printslope@px}
    \edef\nero@printslope@ry{\nero@printslope@ypos}
    \foreach \nero@printslope@order in {#1}{
        % qy = (px/xpos)^order * ypos > ypos  (cote gauche monte ; hypo descend vers la droite)
        \pgfmathparse{
            ((\nero@printslope@px/\nero@printslope@xpos)^(\nero@printslope@order))*\nero@printslope@ypos}
        \edef\nero@printslope@qy{\pgfmathresult}
        % cote gauche montant + hypotenuse ; angle droit en (xpos, ypos)
        \edef\nero@aux1{\noexpand\draw[line width=0.6pt]
            (axis cs:\nero@printslope@xpos,\nero@printslope@ypos)
            -- (axis cs:\nero@printslope@xpos,\nero@printslope@qy)
            -- (axis cs:\nero@printslope@px,\nero@printslope@py);}
        \nero@aux1
        % label de pente sur le cote gauche
        \pgfmathparse{10^((ln(\nero@printslope@ry)+ln(\nero@printslope@qy))/(ln(10)*2))}
        \edef\nero@printslope@labelpos{\pgfmathresult}
        \edef\nero@aux2{\noexpand\node[anchor=east] at
            (axis cs:\nero@printslope@xpos,\nero@printslope@labelpos)
            {\noexpand\tiny \nero@printslope@order};}
        \nero@aux2
        \global\edef\nero@printslope@ry{\nero@printslope@qy}
    }
    % base horizontale partagee, au niveau de ypos (angle droit en bas-gauche)
    \edef\nero@aux1{\noexpand\draw[line width=0.6pt]
        (axis cs:\nero@printslope@xpos,\nero@printslope@ypos)
        |- (axis cs:\nero@printslope@px,\nero@printslope@py);}
    \nero@aux1
    % label "1" sous la base
    \pgfmathparse{10^((ln(\nero@printslope@px)+ln(\nero@printslope@xpos))/(ln(10)*2))}
    \edef\nero@printslope@labelpos{\pgfmathresult}
    \edef\nero@aux2{\noexpand\node[anchor=north] at
        (axis cs:\nero@printslope@labelpos,\nero@printslope@ypos)
        {\noexpand\tiny 1};}
    \nero@aux2
}
\makeatother

%%%%%
% Plot slopes according to h
%%%%%

\makeatletter
\newcommand{\printslopeinv}[4]{
    \tikzset{fixed point arithmetic}
    % get arguments
    \def\nero@printslope@orderlist{#1}
    \edef\nero@printslope@xpos{#2}
    \edef\nero@printslope@ypos{#3}
    \edef\nero@printslope@width{#4}
    % get points position
    \pgfmathparse{\nero@printslope@xpos+\nero@printslope@width}
    \edef\nero@printslope@px{\pgfmathresult}
    \edef\nero@printslope@py{\nero@printslope@ypos}
    \edef\nero@printslope@qx{\pgfmathresult}
    \edef\nero@printslope@ry{\nero@printslope@ypos}
    \foreach \nero@printslope@order in {#1}{
            \pgfmathparse{
                ((\nero@printslope@px/\nero@printslope@xpos)^(\nero@printslope@order))*\nero@printslope@ypos}
            \edef\nero@printslope@qy{\pgfmathresult}
            \edef\nero@aux1{\noexpand\draw[line width=0.6pt]
                (axis cs:\nero@printslope@xpos,\nero@printslope@ypos)
                -- (axis cs:\nero@printslope@qx,\nero@printslope@qy)
                -- (axis cs:\nero@printslope@px,\nero@printslope@py);}
            \nero@aux1
            % slope label
            \pgfmathparse{10^((ln(\nero@printslope@ry)+ln(\nero@printslope@qy))/(ln(10)*2))}
            \edef\nero@printslope@labelpos{\pgfmathresult}
            \edef\nero@aux2{\noexpand\node[anchor=west] at
                (axis cs:\nero@printslope@qx,\nero@printslope@labelpos)
                {\noexpand\tiny \nero@printslope@order};}
            \nero@aux2
            \global\edef\nero@printslope@ry{\nero@printslope@qy}
        }
    % base line
    \edef\nero@aux1{\noexpand\draw[line width=0.6pt]
        (axis cs:\nero@printslope@xpos,\nero@printslope@ypos)
        |- (axis cs:\nero@printslope@px,\nero@printslope@py);}
    \nero@aux1
    % label of base line
    \pgfmathparse{10^((ln(\nero@printslope@px)+ln(\nero@printslope@xpos))/(ln(10)*2))}
    \edef\nero@printslope@labelpos{\pgfmathresult}
    \edef\nero@aux2{\noexpand\node[anchor=north] at
        (axis cs:\nero@printslope@labelpos,\nero@printslope@ypos)
        {\noexpand\tiny 1};}
    \nero@aux2
}
\makeatother

\usepackage{amssymb}
\usepackage{mathtools}
\usepackage{pgfplots}
\usepackage{pgfplotstable}
\usepackage{datatool}
\usepackage{fp}

\pgfplotsset{
    compat=newest,
    smaller labels/.style={
        label style={font=\footnotesize},
        tick label style={font=\footnotesize}
    }
}

\tikzset{font=\small}
\usetikzlibrary{
    fpu,
    fixedpointarithmetic,
    babel,
    external,
    arrows.meta,
    plotmarks,
    positioning,
    angles,
    quotes,
    intersections,
    calc,
    spy,
    decorations.pathreplacing,
    matrix,
    fit,
}
\usepgfplotslibrary{fillbetween}

\gdef\iterator{0}

%%%%%%%%%%%%%%%%%%%%%%%%%%%%%%%%%%%%%
% Plot styles for convergence plots %
%%%%%%%%%%%%%%%%%%%%%%%%%%%%%%%%%%%%%

% Define colors
\definecolor{femcolor}{RGB}{51, 138, 55} %Green (27,158,119)
\definecolor{addcolor}{RGB}{217,95,2} %Orange
\definecolor{add2color}{RGB}{199,39,34} %Red (other prior)
\definecolor{add3color}{RGB}{228, 155, 15} %Dark Yellow (another prior)
\definecolor{multcolor3}{RGB}{117,112,179} %Purple 
\definecolor{multcolor100}{RGB}{0,0,0} %Black (+ empty marker)
\definecolor{multcolor0weak}{RGB}{49, 73, 181} %Blue
\definecolor{multcolor0strong}{RGB}{49, 181, 161} %Cyan

% Define line styles according to the method 
% FEM : solid / Add : dashed / Mult : dotted

% Define marker styles according to the degree
% P1 : square / P2 : circle / P3 : triangle

% Define plot sizes
\newlength{\plotwidth}
\setlength{\plotwidth}{0.45\textwidth}
\newlength{\plotheight}
\setlength{\plotheight}{0.3\textwidth}

%%%%%%%%%%%%%%%%%%%%%%%%%%%%%
% General cvg h environment %
%%%%%%%%%%%%%%%%%%%%%%%%%%%%%

\newenvironment{cvgh}[6][]{
    \begin{tikzpicture}[every node/.style={inner sep=0.5,outer sep=0.5}]
        \edef\filename{#2}
        \edef\legendcolumns{#3}
        \edef\slopes{#4}
        \edef\ypos{#5}
        \edef\normtype{#6} % L2 or H1
        \edef\extraoptions{#1}

        % Read the CSV file into a table
        \pgfplotstableread[col sep=comma]{\filename}\datatable

        % Obtenir le second élément
        \pgfmathtruncatemacro{\secondrow}{1} % Index de la dernière ligne
        \pgfplotstablegetelem{\secondrow}{h}\of\datatable
        \pgfmathsetmacro{\second}{\pgfplotsretval} % Dernière valeur de h_rounded

        % Obtenir le premier élément
        \pgfmathtruncatemacro{\firstrow}{0} % Index de l'avant-dernière ligne
        \pgfplotstablegetelem{\firstrow}{h}\of\datatable
        \pgfmathsetmacro{\first}{\pgfplotsretval} % Avant-dernière valeur de h_rounded

        % Calculer la différence entre les deux
        \pgfmathsetmacro{\diff}{\first - \second}

        %update iterator
        \pgfmathtruncatemacro{\iterator}{\iterator+1}

        \begin{loglogaxis}[
            smaller labels,
            name = left_plot,
            % axis lines
            axis lines = left,
            enlarge x limits={abs=10pt},
            enlarge y limits={abs=10pt},
            % axis x line shift = -5pt,
            % axis y line shift = -5pt,
            % labels
			xmode=log,
            xlabel = {$h$},
            ylabel = {\rotatebox{270}{\normtype}},
            xlabel style={at={(ticklabel* cs:1.01)},anchor=south west},
            ylabel style={at={(ticklabel* cs:1.01)},anchor=west},
            % ticks and labels
            xtick=data,
            xticklabels from table={\datatable}{h},
            % size
            width=\plotwidth, height=\plotheight,
            % marks
            mark options={solid, scale=1},
            % grid
            grid = major,
            % legend
            legend columns=\legendcolumns,
            legend to name=leg:legendFEMCORR_\iterator,
            legend image post style={mark options={solid, scale=1}},
            \extraoptions % <--- Ajout des options supplémentaires ici
        ]
        \expandafter\printslopeinv\expandafter{\slopes}{\second}{\ypos}{\diff}
    }
    {
        \end{loglogaxis}
        \node[yshift=-20pt] at (left_plot.outer south) {\pgfplotslegendfromname{leg:legendFEMCORR_\iterator}};

    \end{tikzpicture}
}

%%%%%%%%%%%%%%%%%%%%%%%%%%%
% FEM vs Add, all degrees %
% L2 norm or H1 semi-norm %
%%%%%%%%%%%%%%%%%%%%%%%%%%%
% use in : ell2D.tex + lap2D_mixed.tex + lap2D.tex (Low frequency + High frequency)

\newcommand{\cvgFEMCorrAlldegGeneral}[5]{
    \edef\fem{#1}
    \edef\add{#2}

    \begin{cvgh}{\fem}{3}{#3}{#4}{#5}
        % Complete the legend
        \addlegendentry{\,FEM $\mathbb{P}_1$\;}
        \addlegendentry{\,FEM $\mathbb{P}_2$\;}
        \addlegendentry{\,FEM $\mathbb{P}_3$\;}
        \addlegendentry{\,Add $\mathbb{P}_1$\;}
        \addlegendentry{\,Add $\mathbb{P}_2$\;}
        \addlegendentry{\,Add $\mathbb{P}_3$\;}

        % Plot FEM
        \addplot [style={solid}, mark=square*, mark size=2, color=femcolor, line width=0.8pt ]
        table [x=h, y=P1, col sep=comma]
            {\fem};
        
        \addplot [style={solid}, mark=*, mark size=2, color=femcolor, line width=0.8pt ]
        table [x=h, y=P2, col sep=comma]
            {\fem};
        
        \addplot [style={solid}, mark=triangle*, mark size=2, color=femcolor, line width=0.8pt ]
        table [x=h, y=P3, col sep=comma]
            {\fem};

        % Plot Add
        \addplot [style={dashed}, mark=square*, mark size=2, color=addcolor, line width=0.8pt ]
        table [x=h, y=P1, col sep=comma]
            {\add};

        \addplot [style={dashed}, mark=*, mark size=2, color=addcolor, line width=0.8pt ]
        table [x=h, y=P2, col sep=comma]
            {\add};

        \addplot [style={dashed}, mark=triangle*, mark size=2, color=addcolor, line width=0.8pt ]
        table [x=h, y=P3, col sep=comma]
            {\add};

    \end{cvgh}
}

% L2 norm
\newcommand{\cvgFEMCorrAlldeg}[3]{
    \cvgFEMCorrAlldegGeneral{#1}{#2}{2,3,4}{#3}{$L^2$}
}

% H1 semi-norm
\newcommand{\cvgFEMCorrAlldegHun}[3]{
    \cvgFEMCorrAlldegGeneral{#1}{#2}{1,2,3}{#3}{Semi-$H^1$}
}